\documentclass[aspectratio=169]{beamer}

\usetheme{default}
\usecolortheme{seahorse}
\setbeamertemplate{navigation symbols}{}
\setbeamertemplate{footline}[frame number]
\setbeamertemplate{itemize items}[circle]

\usepackage[utf8]{inputenc}
\usepackage{amsmath, amssymb, amsthm}
\usepackage{tikz}

\usepackage{graphicx,listings,amsmath,amssymb,amsthm,amsfonts,tikz,tikz-cd,float}
\usetikzlibrary{arrows.meta, positionand, calc, babel}
% General Tikz Settings
\tikzset{
    base node/.style={circle, draw, minimum size=0.85cm,
    font=\sffamily\bfseries\normalsize, thick},
    derived node/.style={circle, draw, minimum size=0.72cm, inner
    sep=0.75pt, font=\sffamily\small\bfseries, thick, fill=white},
    edge label/.style={font=\sffamily\small, fill=white, inner sep=1.25pt},
}

% Colors used consistently for voltage 0 / 1 / 2 throughout this section
\colorlet{volt0}{blue}
\colorlet{volt1}{red}
\colorlet{volt2}{green!55!black}

% Draw the 3 base nodes on a small equilateral triangle
\newcommand{\drawbase}{
    \node[base node] (u) at (0,0) {\( u \)};
    \node[base node] (v) at (3.4,0) {\( v \)};
    \node[base node] (w) at (1.7,2.94) {\( w \)};
}

% Step 1: invisible coordinates for the 3x3 grid, spaced out to
% reduce edge crowding.
% Edges are drawn against these bare coordinates; the visible node
% circles are drawn
% afterwards (see \drawfibernodes) so that edges never paint over a
% node's fill or label.
\newcommand{\drawfibers}{
    \foreach \name/\x in {u/0, v/3.8, w/7.6} {
        \coordinate (\name0) at (\x, 5.1);
        \coordinate (\name1) at (\x, 2.55);
        \coordinate (\name2) at (\x, 0);
    }
}

% Step 2: the visible fiber labels and node circles, drawn last so
% they sit on top of
% every edge.
\newcommand{\drawfibernodes}{
    \foreach \name/\x in {u/0, v/3.8, w/7.6} {
        \node[font=\sffamily\bfseries\small, text=gray] at (\x, 6.2)
        {Fiber \name};
        \node[derived node, draw=volt0, text=volt0] at (\name0) {\(\name,0\)};
        \node[derived node, draw=volt1, text=volt1] at (\name1) {\(\name,1\)};
        \node[derived node, draw=volt2, text=volt2] at (\name2) {\(\name,2\)};
    }
}

% Draw the undirected intra-fiber loops (lifts of base loops form a
% C_3 in each fiber)
\newcommand{\drawderivedloops}{
    \foreach \name in {u, v, w} {
        \draw[-, thick, black, opacity=0.85] (\name0) to (\name1);
        \draw[-, thick, black, opacity=0.85] (\name1) to (\name2);
        \draw[-, thick, black, opacity=0.85, bend right=40]
        (\name0) to (\name2);
    }
}

\newcommand{\drawlift}[3]{
    % #1 = source fiber, #2 = target fiber, #3 = voltage
    \ifnum#3=0
    \draw[->, >=stealth, thick, volt0, opacity=0.7,
    shorten >=2.2pt, shorten <=2.2pt]
    (#10) to (#20);
    \draw[->, >=stealth, thick, volt0, opacity=0.7,
    shorten >=2.2pt, shorten <=2.2pt]
    (#11) to (#21);
    \draw[->, >=stealth, thick, volt0, opacity=0.7,
    shorten >=2.2pt, shorten <=2.2pt]
    (#12) to (#22);
    \fi
    \ifnum#3=1
    \draw[->, >=stealth, thick, volt1, opacity=0.6,
    shorten >=2.2pt, shorten <=2.2pt]
    (#10) to (#21);
    \draw[->, >=stealth, thick, volt1, opacity=0.6,
    shorten >=2.2pt, shorten <=2.2pt]
    (#11) to (#22);
    \draw[->, >=stealth, thick, volt1, opacity=0.6,
    shorten >=2.2pt, shorten <=2.2pt]
    (#12) to (#20);
    \fi
    \ifnum#3=2
    \draw[->, >=stealth, thick, volt2, opacity=0.6,
    shorten >=2.2pt, shorten <=2.2pt]
    (#10) to (#22);
    \draw[->, >=stealth, thick, volt2, opacity=0.6,
    shorten >=2.2pt, shorten <=2.2pt]
    (#11) to (#20);
    \draw[->, >=stealth, thick, volt2, opacity=0.6,
    shorten >=2.2pt, shorten <=2.2pt]
    (#12) to (#21);
    \fi
}

\newcommand{\drawliftgrey}[3]{
    % #1 = source fiber, #2 = target fiber, #3 = voltage
    \ifnum#3=0
    \draw[->, >=stealth, thick, darkgrey, opacity=0.7,
    shorten >=2.2pt, shorten <=2.2pt]
    (#10) to (#20);
    \draw[->, >=stealth, thick, darkgrey, opacity=0.7,
    shorten >=2.2pt, shorten <=2.2pt]
    (#11) to (#21);
    \draw[->, >=stealth, thick, darkgrey, opacity=0.7,
    shorten >=2.2pt, shorten <=2.2pt]
    (#12) to (#22);
    \fi
    \ifnum#3=1
    \draw[->, >=stealth, thick, darkgrey, opacity=0.6,
    shorten >=2.2pt, shorten <=2.2pt]
    (#10) to (#21);
    \draw[->, >=stealth, thick, darkgrey, opacity=0.6,
    shorten >=2.2pt, shorten <=2.2pt]
    (#11) to (#22);
    \draw[->, >=stealth, thick, darkgrey, opacity=0.6,
    shorten >=2.2pt, shorten <=2.2pt]
    (#12) to (#20);
    \fi
    \ifnum#3=2
    \draw[->, >=stealth, thick, darkgrey, opacity=0.6,
    shorten >=2.2pt, shorten <=2.2pt]
    (#10) to (#22);
    \draw[->, >=stealth, thick, darkgrey, opacity=0.6,
    shorten >=2.2pt, shorten <=2.2pt]
    (#11) to (#20);
    \draw[->, >=stealth, thick, darkgrey, opacity=0.6,
    shorten >=2.2pt, shorten <=2.2pt]
    (#12) to (#21);
    \fi
}

% One-line color legend, reused above every derived graph
\newcommand{\voltagelegend}{%
    \begin{tikzpicture}[baseline]
        \node[font=\sffamily\small] at (0,0)
        {\textcolor{volt0}{\textbf{---}} voltage 0 \quad
            \textcolor{volt1}{\textbf{---}} voltage 1 \quad
        \textcolor{volt2}{\textbf{---}} voltage 2};
    \end{tikzpicture}%
}


\title{Quantum Graphs over $3\times 3$ Matrices}
\subtitle{Classification, Voltage Graphs, and Comparisons}
\author{Filip Rabiega, Dr. Mariusz Tobolski (supervisor)}
\date{September 11th, 2026}

\begin{document}

% ---------------------------------------------------------
\begin{frame}
    \titlepage
\end{frame}

% ---------------------------------------------------------
\begin{frame}{Roadmap}
    \begin{itemize}
        \item \(C^{*}\)-algebras and crossed product \(C^{*}\)-algebras
        \item Quantum sets and quantum graphs \hfill \emph{(brief recap)}
        \item Gromada's theorem: quantum graphs on $M_n(\mathbb{C})$
            $\leftrightarrow$ subspaces
        \item Voltage and derived quantum graphs
        \item Quantum isomorphism
        \item[]
        \item \textbf{Main content: quantum graphs on \( M_3 (\mathbb{C}) \)}
            \begin{itemize}
                \item Classification via clock--shift matrices
                \item Realization as classical derived graphs
                \item Vertex-transitivity
                \item Where the story breaks down
            \end{itemize}
    \end{itemize}
\end{frame}

% ---------------------------------------------------------
\begin{frame}{\(C^{*}\)-algebras and crossed product \(C^{*}\)-algebras}
    \textbf{\(C^{*}\)-algebra:} a Banach algebra \( A \) satisfying
    the \(C^{*}\)-axiom: \( \left\| a^* a \right\| = \left\| a
    \right\| \) for all \( a \in A \). I only consider unital
    \(C^{*}\)-algebras.

    \vspace{0.5em}
    \textbf{Crossed product \(C^{*}\)-algebra:} a \(C^{*}\)-algebra
    \( A \rtimes_{\alpha} G \) spanned
    by the  functions \( u_g: G \rightarrow A \),
    \[
        u_g(h) =
        \begin{cases}
            1_A, & h = g, \\
            0_A, & h \neq g,
        \end{cases}
        \qquad h \in G,
    \]
    where \( G \) is a finite group. Multiplication, involution and norm are
    determined by an action \( \alpha \) of \( G \) on \( A \).
\end{frame}

% ---------------------------------------------------------
\begin{frame}{Quantum sets and quantum graphs}
    \textbf{Quantum set:} a pair $(B, \psi)$ -- finite-dim.\ $C^*$-algebra $B$
    with a positive and faithful functional $\psi$ such that \( mm* =
    \mathrm{Id} \).
    \begin{itemize}
        \item Classical set $V$ $\rightsquigarrow$ $(\mathbb{C}^V, \psi_V)$
        \item $M_n(\mathbb{C})$ $\rightsquigarrow$ $(M_n(\mathbb{C}),
            n\operatorname{Tr})$ -- a
            \emph{genuinely noncommutative} example
    \end{itemize}
    \vspace{0.5em}
    \textbf{Quantum graph:} a linear map $A \colon B \to B$ generalizing an
    adjacency matrix, with
    \[
        m(A\otimes A)m^* = A, \qquad A(b^*) = A(b)^*.
    \]
    \begin{itemize}
        \item \emph{Loop-free}: $m(A \otimes \mathrm{Id})m^* = 0$
        \item \emph{Undirected}: $A = A^*$
    \end{itemize}
\end{frame}

% ---------------------------------------------------------
\begin{frame}{Gromada's theorem}
    \begin{block}{Theorem (Gromada)}
        Quantum graphs on $M_n(\mathbb{C})$ with $m$ quantum edges
        $\longleftrightarrow$
        $m$-dimensional subspaces $V \subseteq \mathfrak{gl}_n$, via
        \[
            A_V = n^2 \sum_{s=1}^m \xi_s \otimes \bar{\xi_s}.
        \]
        \begin{itemize}
            \item No loops $\iff$ $V \subseteq \mathfrak{sl}(n) =
                \left\{ A \in \mathfrak{gl}(n) : \operatorname{Tr}(A)
                = 0 \right\}$
            \item Undirected $\iff$ $V = V^*$
            \item Isomorphic graphs $\iff$ $V = UWU^*$, $U \in SU(n)
                = \left\{ U \in U(n) : \det(U) = 1 \right\}$
        \end{itemize}
    \end{block}
    \vspace{0.5em}
    This turns questions about quantum graphs into \textbf{linear algebra}.
\end{frame}

% ---------------------------------------------------------
\begin{frame}{Voltage and derived quantum graphs}
    \textbf{Classical picture (Gross--Tucker):} label edges of a multigraph
    $\Gamma$ by a group $M_3 (\mathbb{C})$ $\rightsquigarrow$ unfold
    into a derived graph
    $\Gamma^\lambda$ on $V \times G$.

    \vspace{0.5em}
    \textbf{Quantum generalization:} a $M_3 (\mathbb{C})$-indexed family
    $\tilde{A} = (\tilde{A}_g)_{g\in G}$ of quantum adjacency matrices,
    reassembled via a crossed product
    \[
        \tilde{A} \rtimes_{\hat\alpha} \hat{G} := \sum_{g \in G}
        m(X_g \otimes E_\rtimes^* \tilde{A}_g E_\rtimes)\, m^*.
    \]

    \vspace{0.5em}
    \begin{block}{Key result (Schäfer--Tobolski)}
        Derived quantum graphs w.r.t.\ \emph{different} actions on the same
        voltage data are always \textbf{quantum isomorphic}.
    \end{block}
\end{frame}

% ---------------------------------------------------------
\begin{frame}{Example of a voltage and a derived graph}
    \begin{figure}[H]
        % \centering
        % {\normalsize\sffamily\bfseries Directed base graph over \(
        % \mathbb{Z}_3 \) (single vertex, two loops)}\\[4pt]
        \begin{minipage}[t]{0.435\textwidth}
            \centering
            \begin{tikzpicture}[>=Stealth, scale=1.0]
                \node[base node] (v) at (0,0) {\( v \)};
                \path[->, thick]
                (v) edge [loop left, distance=1.7cm, in=145, out=215]
                node[edge label] {\( 1 \)} (v)
                (v) edge [loop right, distance=1.7cm, in=35, out=-35]
                node[edge label] {\( 2 \)} (v);
            \end{tikzpicture}\\[6pt]
            \captionof{}{Base voltage graph \( \Gamma \) over \(
            \mathbb{Z}_3 \)}
        \end{minipage}\hfill
        \begin{minipage}[t]{0.515\textwidth}
            \centering
            \begin{tikzpicture}[>=Stealth, scale=0.72]
                \node[derived node, draw=volt0, text=volt0] (v0) at
                (90:2.6)  {\( v,0 \)};
                \node[derived node, draw=volt1, text=volt1] (v1) at
                (210:2.6) {\( v,1 \)};
                \node[derived node, draw=volt2, text=volt2] (v2) at
                (330:2.6) {\( v,2 \)};

                % e1 lifts (voltage 1): v0 -> v1 -> v2 -> v0
                \path[->, thick, volt1, opacity=0.6]
                (v0) edge [bend right=20] (v1)
                (v1) edge [bend right=20] (v2)
                (v2) edge [bend right=20] (v0);

                % e2 lifts (voltage 2): v0 -> v2 -> v1 -> v0
                \path[->, thick, volt2, opacity=0.6]
                (v0) edge [bend right=20] (v2)
                (v2) edge [bend right=20] (v1)
                (v1) edge [bend right=20] (v0);
            \end{tikzpicture}\\[6pt]
            \captionof{}{Derived graph \( \Gamma^\lambda \) (two \(
            3 \)-cycles)}
        \end{minipage}
        \caption{Voltage graph and derived cover over \( \mathbb{Z}_3 \),
        for a base graph with a single vertex and two loops.}
        \label{fig:voltage-graph-z3-loops}
    \end{figure}
\end{frame}

% ---------------------------------------------------------
\begin{frame}{Quantum isomorphism}
    Two quantum graphs $A_1, A_2$ are \textbf{quantum isomorphic},
    $A_1 \cong_q A_2$, if there is a unitary
    $\mathcal{H}$-twisted map $\rho$ intertwining them.
    \begin{itemize}
        \item Taking $\mathcal{H} = \mathbb{C}$ recovers ordinary isomorphism
        \item $\cong_q$ is \textbf{strictly weaker} than $\cong$
        \item Originates in nonlocal games: quantum strategies can win games
            classical strategies cannot
    \end{itemize}
    \vspace{0.5em}
    This is the tool that lets us relate ``truly quantum'' graphs on
    $M_3 (\mathbb{C})$
    to \emph{classical} graphs.
\end{frame}

% ---------------------------------------------------------
\begin{frame}{Generalised Pauli matrices on $M_3 (\mathbb{C})$}
    Clock and shift matrices:
    \[
        P =
        \begin{pmatrix} 1 & 0 & 0 \\ 0 & \omega & 0 \\ 0 & 0 & \omega^2
        \end{pmatrix},
        \qquad
        Q =
        \begin{pmatrix} 0 & 0 & 1 \\ 1 & 0 & 0 \\ 0 & 1 & 0
        \end{pmatrix},
        \qquad \omega = e^{2\pi i/3}.
    \]
    \begin{itemize}
        \item $P^3 = Q^3 = I$, \quad $PQ = \omega\, QP$
        \item $\{P^aQ^b\}_{a,b\in\{0,1,2\}}$ is an \emph{orthogonal basis}
            of $M_3 (\mathbb{C})$
        \item Higher-dimensional analogue of the Pauli matrices
    \end{itemize}
\end{frame}

% ---------------------------------------------------------
\begin{frame}{Classifying subspaces generated by $P, Q$}
    Reflection-invariant blocks: $\mathcal{P}_1 = \{P,P^2\}$,
    $\mathcal{P}_2 = \{Q,Q^2\}$, $\mathcal{P}_3, \mathcal{P}_4$ (mixed).

    \begin{center}
        \begin{tabular}{ccc}
            \toprule
            $m$ & subspaces & count \\
            \midrule
            $0$ & $V_0 = \{0\}$ & $1$ \\
            $2$ & $V_k = \operatorname{span}(\mathcal{P}_k)$ & $4$ \\
            $4$ & $W_{i,j} = \operatorname{span}(\mathcal{P}_i \oplus
            \mathcal{P}_j)$ & $6$ \\
            $6$ & $U_k = \operatorname{span}(\cup_{i \neq k}
            \mathcal{P}_i)$ & $4$ \\
            $8$ & $\mathfrak{sl}_3$ & $1$ \\
            \bottomrule
        \end{tabular}
    \end{center}
    \vspace{0.3em}
    $16$ self-reflective subspaces total, built from $2^4$ block choices.
\end{frame}

% ---------------------------------------------------------
\begin{frame}{Full unitary equivalence}
    The discrete Fourier transform
    \[
        F = \frac{1}{\sqrt{3}}
        \begin{pmatrix} 1 & 1 & 1 \\ 1 & \omega & \omega^2 \\ 1 &
            \omega^2 & \omega
        \end{pmatrix}
    \]
    is not enough.
    \vspace{0.3em}
    \begin{block}{Key trick}
        A second unitary $S$ (a phase gate) cyclically permutes
        $\mathcal{P}_2, \mathcal{P}_3, \mathcal{P}_4$.
        Together, $F$ and $S$ generate the full alternating group $A_4$
    \end{block}

    \vspace{0.3em}
    \begin{block}{Corollary}
        All 16 subspaces collapse to exactly \textbf{5} isomorphism classes
        -- one per dimension $m \in \{0,2,4,6,8\}$.
    \end{block}
\end{frame}

% ---------------------------------------------------------
\begin{frame}{Realizing them as classical graphs}
    Each nontrivial class ($m=2,4,6,8$) is exhibited as the \textbf{derived
    quantum graph} of an explicit classical voltage graph over $\mathbb{Z}_3$.

    \begin{itemize}
        \item $m=2$: loops only
        \item $m=4$: loops + oppositely-oriented triangles
        \item $m=6$: loops + partial $K_{3,3}$ between fibers
        \item $m=8$: full $K_{3,3}$ (complete graph $\mathfrak{sl}_3$)
    \end{itemize}

    \vspace{0.5em}
    By a theorem of Schäfer and Tobolski: each is \textbf{quantum isomorphic}
    to a classical graph on \textbf{9 vertices}.
\end{frame}

% ---------------------------------------------------------
\begin{frame}{Classical graphs quantum isomorphic to our graphs}
    % ==========================================
    % CASE 1: Only 1- and 2- loops
    % ==========================================
    \begin{figure}[H]
        \begin{minipage}[c]{0.435\textwidth}
            \centering
            \begin{tikzpicture}[scale=0.85]
                \drawbase
                \path[->, >=stealth, thick]
                (u) edge [loop left, looseness=8] node[edge label]
                {\( 1,2 \)} (u)
                (v) edge [loop right, looseness=8] node[edge label]
                {\( 1,2 \)} (v)
                (w) edge [loop above, looseness=8] node[edge label]
                {\( 1,2 \)} (w);
            \end{tikzpicture}
            \captionof{}{Base \( \Gamma \): directed loops}
        \end{minipage}\hfill
        \begin{minipage}[c]{0.515\textwidth}
            \centering
            \begin{tikzpicture}[scale=0.62]
                \drawfibers
                \drawderivedloops
                \drawfibernodes
            \end{tikzpicture}
            \captionof{}{Derived \( \Gamma^\lambda \): undirected 3-cycles}
        \end{minipage}
        \caption{Voltage graph and derived cover for the loopfree,
            undirected quantum graph on \( M_3 (\mathbb{C}) \) with \( m=2 \)
        quantum edges.}
        \label{fig:voltage-graph-m2-edges}
    \end{figure}
\end{frame}

% ---------------------------------------------------------
\begin{frame}
    % ==========================================
    % CASE 2: 1- and 2- cycles along with loops
    % ==========================================
    \begin{figure}[H]
        \begin{minipage}[c]{0.435\textwidth}
            \centering
            \begin{tikzpicture}[scale=0.85]
                \drawbase
                \path[->, >=stealth, thick]
                (u) edge [loop left, looseness=8] node[edge label]
                {\( 1,2 \)} (u)
                (v) edge [loop right, looseness=8] node[edge label]
                {\( 1,2 \)} (v)
                (w) edge [loop above, looseness=8] node[edge label]
                {\( 1,2 \)} (w)
                (u) edge [bend left=25] node[edge label] {\( 1 \)} (v)
                (v) edge [bend left=25] node[edge label] {\( 2 \)} (u)
                (v) edge [bend left=25] node[edge label] {\( 1 \)} (w)
                (w) edge [bend left=25] node[edge label] {\( 2 \)} (v)
                (w) edge [bend left=25] node[edge label] {\( 1 \)} (u)
                (u) edge [bend left=25] node[edge label] {\( 2 \)} (w);
            \end{tikzpicture}
            \captionof{}{Base \( \Gamma \): directed triangle \& loops}
        \end{minipage}\hfill
        \begin{minipage}[c]{0.515\textwidth}
            \centering
            \begin{tikzpicture}[scale=0.62]
                \drawfibers
                \drawderivedloops
                \drawliftgrey{u}{v}{1}
                \drawliftgrey{v}{w}{1}
                \drawliftgrey{w}{u}{1}
                \drawliftgrey{u}{w}{2}
                \drawliftgrey{v}{u}{2}
                \drawliftgrey{w}{v}{2}
                \drawfibernodes
            \end{tikzpicture}
            \captionof{}{Derived \( \Gamma^\lambda \)}
        \end{minipage}
        \caption{Voltage graph and derived cover for the loopfree,
            undirected quantum graph on \( M_3 (\mathbb{C}) \) with \( m=4 \)
        quantum edges.}
        \label{fig:voltage-graph-m4-edges}
    \end{figure}
\end{frame}

% ---------------------------------------------------------
\begin{frame}
    % ==========================================
    % CASE 3: All 1's and all 2's
    % ==========================================
    \begin{figure}[H]
        \begin{minipage}[c]{0.435\textwidth}
            \centering
            \begin{tikzpicture}[scale=0.85]
                \drawbase
                \path[thick]
                (u) edge [bend left=25] node[edge label] {\( 1 \)} (v)
                (u) edge [bend right=25] node[edge label] {\( 2 \)} (v)
                (v) edge [bend left=25] node[edge label] {\( 1 \)} (w)
                (v) edge [bend right=25] node[edge label] {\( 2 \)} (w)
                (w) edge [bend left=25] node[edge label] {\( 1 \)} (u)
                (w) edge [bend right=25] node[edge label] {\( 2 \)} (u);
                \path[->, >=stealth, thick]
                (u) edge [loop left, looseness=8] node[edge label]
                {\( 1,2 \)} (u)
                (v) edge [loop right, looseness=8] node[edge label]
                {\( 1,2 \)} (v)
                (w) edge [loop above, looseness=8] node[edge label]
                {\( 1,2 \)} (w)
            \end{tikzpicture}
            \captionof{}{Base \( \Gamma \): directed multigraph}
        \end{minipage}\hfill
        \begin{minipage}[c]{0.515\textwidth}
            \centering
            \begin{tikzpicture}[scale=0.62]
                \drawfibers
                \drawderivedloops
                \drawliftgrey{u}{v}{1}
                \drawliftgrey{u}{w}{1}
                \drawliftgrey{v}{u}{1}
                \drawliftgrey{v}{w}{1}
                \drawliftgrey{w}{u}{1}
                \drawliftgrey{w}{v}{1}
                \drawliftgrey{u}{v}{2}
                \drawliftgrey{u}{w}{2}
                \drawliftgrey{v}{u}{2}
                \drawliftgrey{v}{w}{2}
                \drawliftgrey{w}{u}{2}
                \drawliftgrey{w}{v}{2}
                % \foreach \src/\dst in {u/v, v/w, w/u} {
                %     \drawlift{\src}{\dst}{1}
                %     \drawlift{\src}{\dst}{2}
                % }
                \drawfibernodes
            \end{tikzpicture}
            \captionof{}{Derived \( \Gamma^\lambda \): partial
            \(K_{3,3}\) between fibers}
        \end{minipage}
        \caption{Voltage graph and derived cover for the loopfree,
            undirected quantum graph on \( M_3 (\mathbb{C}) \) with \( m=6 \)
        quantum edges.}
        \label{fig:voltage-graph-m6-edges}
    \end{figure}
\end{frame}

% ---------------------------------------------------------
\begin{frame}
    % ==========================================
    % CASE 4: Everything except 0-loops
    % ==========================================
    \begin{figure}[H]
        \begin{minipage}[c]{0.435\textwidth}
            \centering
            \begin{tikzpicture}[scale=0.85]
                \drawbase
                \path[thick]
                (u) edge [bend left=32] node[edge label] {\( 1 \)} (v)
                (u) edge node[edge label] {\( 0 \)} (v)
                (u) edge [bend right=32] node[edge label] {\( 2 \)} (v)
                (v) edge [bend left=32] node[edge label] {\( 1 \)} (w)
                (v) edge node[edge label] {\( 0 \)} (w)
                (v) edge [bend right=32] node[edge label] {\( 2 \)} (w)
                (w) edge [bend left=32] node[edge label] {\( 1 \)} (u)
                (w) edge node[edge label] {\( 0 \)} (u)
                (w) edge [bend right=32] node[edge label] {\( 2 \)} (u);
                \path[->, >=stealth, thick]
                (u) edge [loop left, looseness=8] node[edge label]
                {\( 1,2 \)} (u)
                (v) edge [loop right, looseness=8] node[edge label]
                {\( 1,2 \)} (v)
                (w) edge [loop above, looseness=8] node[edge label]
                {\( 1,2 \)} (w)
            \end{tikzpicture}
            \captionof{}{Base \( \Gamma \): full directed multigraph}
        \end{minipage}\hfill
        \begin{minipage}[c]{0.515\textwidth}
            \centering
            \begin{tikzpicture}[scale=0.62]
                \drawfibers
                \drawderivedloops
                \drawliftgrey{u}{v}{0}
                \drawliftgrey{u}{w}{0}
                \drawliftgrey{v}{u}{0}
                \drawliftgrey{v}{w}{0}
                \drawliftgrey{w}{u}{0}
                \drawliftgrey{w}{v}{0}
                \drawliftgrey{u}{v}{1}
                \drawliftgrey{u}{w}{1}
                \drawliftgrey{v}{u}{1}
                \drawliftgrey{v}{w}{1}
                \drawliftgrey{w}{u}{1}
                \drawliftgrey{w}{v}{1}
                \drawliftgrey{u}{v}{2}
                \drawliftgrey{u}{w}{2}
                \drawliftgrey{v}{u}{2}
                \drawliftgrey{v}{w}{2}
                \drawliftgrey{w}{u}{2}
                \drawliftgrey{w}{v}{2}
                % \foreach \src/\dst in {u/v, v/w, w/u} {
                %     \drawlift{\src}{\dst}{0}
                %     \drawlift{\src}{\dst}{1}
                %     \drawlift{\src}{\dst}{2}
                % }
                \drawfibernodes
            \end{tikzpicture}
            \captionof{}{Derived \( \Gamma^\lambda \): perfect
            \(K_{3,3}\) between all fibers}
        \end{minipage}
        \caption{Voltage graph and derived cover for the loopfree,
            undirected quantum graph on \( M_3 (\mathbb{C}) \) with \( m=8 \)
        quantum edges.}
        \label{fig:voltage-graph-m8-edges}
    \end{figure}
\end{frame}

% ---------------------------------------------------------
\begin{frame}{Vertex-transitivity}
    \textbf{Vertex-transitive quantum graph:} a quantum graph whose Gromada
    subspace \( V \) is such that \( \operatorname{span}(\left\{ u
    \in U(n) : uVu^* = V \right\}) = M_3 (\mathbb{C})\).
    \begin{block}{Proposition}
        Every self-reflective subspace generated by $P, Q$ gives a
        \textbf{vertex-transitive} quantum graph.
    \end{block}
    \begin{itemize}
        \item Proof idea: conjugation by any $P^aQ^b$ fixes each generator
            up to phase $\Rightarrow$ the 9-element group
            $\{P^aQ^b\}$ stabilizes \emph{every} such subspace
        \item Holds already at minimal dimension $m=2$
    \end{itemize}
    \vspace{0.5em}
    \textbf{Comparison:} Hayes et al. classify \emph{all}
    vertex-transitive quantum graphs on $M_3 (\mathbb{C})$ -- a
    strictly larger family
    than ours.
\end{frame}

% ---------------------------------------------------------
\begin{frame}{Where the $M_2(\mathbb{C})$ story breaks down}
    \textbf{$M_2(\mathbb{C})$:} simple graphs classified purely by $m$
    (via $PU(2) \cong SO(3)$).

    \vspace{0.5em}
    \begin{block}{Lemma}
        $PU(3) \curvearrowright \mathbb{P}(\mathfrak{su}(3))$ has
        \textbf{infinitely many orbits}.
    \end{block}
    \begin{itemize}
        \item Proved via a projective invariant $I(X) = c(X)^2/q(X)^3$
        \item Consequence: infinite families of mutually non-isomorphic
            undirected graphs on $M_3 (\mathbb{C})$ for fixed $m$
    \end{itemize}
\end{frame}

% ---------------------------------------------------------
\begin{frame}{Summary}
    \begin{itemize}
        \item Classified $P,Q$-generated quantum graphs on $M_3
            (\mathbb{C})$: exactly
            \textbf{5} isomorphism classes
        \item Realized them as \textbf{classical} graphs on 9 vertices,
            via voltage/derived quantum graphs
        \item All are \textbf{vertex-transitive}
        \item Showed the clean $M_2(\mathbb{C})$ picture
            \textbf{fails} for $M_3 (\mathbb{C})$:
            infinite families exist, and non-classical examples exist
    \end{itemize}
    \vspace{0.7em}
    \textbf{Further work:} extend to $M_n(\mathbb{C})$, $n \geq 4$ --
    no exceptional
    isomorphism like $PU(2)\cong SO(3)$ is available.
\end{frame}

% ---------------------------------------------------------
\begin{frame}
    \centering
    \Large Thank you!
\end{frame}

\end{document}
